% section:question
The question of this physics edition is which particle and geometric readouts
can be consistently constructed from a common graded carrier, and which
additional physical assumptions remain necessary. A common state space,
a mathematical symmetry and a physical identification are separate claims.
This edition is a research draft. The long edition retains the complete
manuscript, alternative routes and historical qualifications.
% section:core
On a declared real Hilbert space let $J$ be an orthogonal involution and
put $C=(I+J)/2$ and $P_-=(I-J)/2$. Then
\[
C^2=(I+2J+J^2)/4=C,\qquad C^*=C,\qquad CP_-=0.
\]
These are algebraic projection identities. Identifying either range with
a physical particle sector requires an additional interpretation. The
generation triple must not be identified with the matter/light/tachyon
sector triple merely because both have three entries. Frame comparisons
require an explicit map and a declared metric.
% section:method
QAM supplies the object and assertion language used by the project. MML
specifies typed instructions, and MMA supplies their bounded execution
semantics and receipts. The local statement store organizes versions,
sources and dependencies; it is not the whole QAM language and is not
a general proof assistant. A source hash fixes content, not truth.

For the projection identity above, record the domain, the assumptions
$J^2=I$ and $J^*=J$, the definition of $C$, and the asserted conclusion.
The displayed calculation is its mathematical argument. A concrete
rational matrix check can independently verify a finite instance: the
registered matrices below check a projector and compare commuting and
noncommuting transports. Those receipts establish only their exact matrix
equations, not a native field realization or a mass mechanism.

In the store, \texttt{requires} records a prerequisite and
\texttt{proves} associates proof text. The interpretation relation names
a declared definition; its store label is \texttt{interpreted-by}.
These are descriptions of store relations, not
new MML instruction spellings. A guard states a missing obligation.
Source-matching, self-review, exact computation and physical interpretation
are different evidence levels. Unreviewed manuscript imports remain
candidate material even when their original headings say ``theorem''.
% section:particles
The particle branch proposes algebraic gauge and generation constructions
on the common carrier. Their chosen embeddings, scalar field and generation
dictionary must be specified. The statements selected for this section
are registered below as candidates pending interpretation, proof and
dependency review; the inventory does not itself establish a derivation
of the Standard Model or a unique number of generations.
% section:masses
Mass readout is a concrete test of the interface between algebra and
observables. The following source-pinned Koide and HC arguments separate
equal-norm identities, conditional root representations, positivity and
the missing native coupling. Their self-reviews apply under the stated
mathematical hypotheses. A fitted root operator or inverse reconstruction
does not determine its physical origin, phase or absolute mass scale.
% section:geometry
The dynamics and geometry branch studies projected evolution and
conditional local field comparisons. A formal master equation does not
by itself specify a self-adjoint Hamiltonian, a well-posed domain or a
regional solution. Selected autonomy and Cartan statements are retained
as candidates until their exact assumptions, proof associations and
resources are mapped into this edition.
% section:limits
The native carrier action, positive three-generation root readout,
phase/scale selection and the required frame bridges remain open.
The recent zeta and uncertainty constructions are comparison models,
not a recovered lepton spectrum or a proof of the historical zeta-clock
claims. Cosmological, Hawking and string/LQG interpretations are retained
in the long edition with their source qualifications. No OP closure is
created by assigning statements to a document or by executing a check.
